\chapter{Telemetria hardware: il modulo \texttt{Sia.HealthChecks.Servers}}
\label{cap:telemetria-hardware}

\intro{Il presente capitolo descrive l'applicazione concreta delle convenzioni di istrumentazione descritte nei capitoli precedenti al modulo \texttt{Sia.HealthChecks.Servers}, dedicato al monitoraggio delle risorse hardware dei server di installazione (CPU, temperatura, RAM, spazio su disco, connessioni SQL). Vengono illustrate le scelte di progetto specifiche per la telemetria hardware, i diagrammi delle classi coinvolte e la dashboard Grafana risultante.}\\

\section{Contesto e obiettivo}

Il modulo \texttt{Sia.HealthChecks.Servers} è parte del framework \emph{Sia.Server.Framework}
ed è responsabile della raccolta periodica di indicatori sullo stato del server che ospita
l'applicazione, con successiva notifica ad un \emph{manager} centrale delle installazioni.
Prima dell'intervento di telemetria, il modulo esponeva già i dati tramite:
\begin{itemize}
    \item cinque implementazioni di \texttt{IHealthCheck} (\texttt{CpuUsageHealthCheck},
        \texttt{CpuTemperatureHealthCheck}, \texttt{RamHealthCheck}, \texttt{DiskHealthCheck},
        \texttt{SqlHealthCheck}), invocate dall'endpoint \texttt{/health} di ASP.NET Core;
    \item un \texttt{BackgroundService} (\texttt{HealthCheckServerHostedService}) che, ad
        intervallo configurabile (\texttt{HealthNotifyDelayMinutes}), aggrega i risultati e
        invia il report al \emph{manager} tramite un \texttt{HttpClient} dedicato
        (\texttt{HealthCheckServerToServerApi}).
\end{itemize}

L'obiettivo dell'istrumentazione è trasformare questi dati in segnali di osservabilità
standard (\emph{trace}, metriche, \emph{log}), permettendone la visualizzazione unificata
su Grafana e la correlazione con le altre trace applicative tramite \texttt{TraceId}.

\section{Scelte specifiche per la telemetria hardware}

Il modulo presenta caratteristiche peculiari che hanno richiesto scelte progettuali diverse
rispetto agli altri moduli instrumentati:

\begin{description}
    \item[Assenza di controller.] Il modulo non espone un proprio \texttt{BaseApiController}:
        gli \texttt{IHealthCheck} sono invocati dall'infrastruttura ASP.NET Core e
        l'\texttt{HttpClient} produce trace automatiche. Di conseguenza, come previsto
        dalla regola del framework, la classe \texttt{SiaHealthChecksServersDiagnostics} non
        implementa l'interfaccia \texttt{IControllerDiagnostics} e non espone il
        \emph{counter} \texttt{ControllerErrors}.
    \item[Valori di stato e non cumulativi.] La maggior parte delle informazioni esposte
        (CPU\%, temperatura, GB di RAM liberi, GB su disco, connessione SQL \emph{up}/\emph{down})
        sono grandezze \emph{istantanee}, non cumulative. L'uso di \texttt{Counter} o
        \texttt{Histogram} sarebbe improprio: serve lo strumento \texttt{ObservableGauge}.
    \item[Aggiornamento on-demand.] Le letture hardware avvengono solo quando l'endpoint
        \texttt{/health} viene invocato (probe esterno, polling del \emph{manager}). Non
        esiste un ciclo interno che campioni i sensori a intervallo fisso. Le
        \texttt{ObservableGauge} devono quindi leggere un valore
        \textbf{cached} in memoria, aggiornato dagli \texttt{IHealthCheck} alla loro
        esecuzione, evitando di esporre valori \emph{stantii} quando nessuna lettura è mai
        stata effettuata.
    \item[Cardinalità controllata.] I \emph{tag} delle metriche hanno domini piccoli e
        chiusi: \texttt{sia.check} $\in\{cpu, cpu\_temperature, ram, disk, sql\}$,
        \texttt{sia.status} $\in\{Healthy, Unhealthy\}$, \texttt{sia.result}
        $\in\{ok, http\_error, exception\}$. I \emph{tag} \texttt{sia.drive} e
        \texttt{sia.connection} sono limitati ai dischi e alle connessioni effettivamente
        configurate (tipicamente $\le 5$).
\end{description}

\section{Segnali di telemetria prodotti}

\subsection{Trace}

Il modulo espone tre \texttt{ActivitySource} custom, tutti dalla stessa classe
\texttt{Sia\-Health\-Checks\-Servers\-Diagnostics}:

\begin{table}[h]
\centering
\caption{Span prodotte dal modulo \texttt{Sia.HealthChecks.Servers}}
\label{tab:span-healthchecks}
\begin{tabularx}{\textwidth}{|l|X|X|}
\hline
\textbf{Span} & \textbf{Contesto} & \textbf{Tag principali} \\
\hline
\texttt{sia.healthchecks.servers.cycle} &
\emph{Root span} di ogni iterazione del \texttt{HostedService} di notifica &
\texttt{sia.operation=notify\_cycle}, \texttt{sia.enabled}, \texttt{sia.delay\_minutes}, \texttt{sia.success}, \texttt{exception.type} \\
\hline
\texttt{sia.healthchecks.servers.check} &
\emph{Child span} per ciascuna invocazione di un \texttt{IHealthCheck} &
\texttt{sia.operation=check}, \texttt{sia.check}, \texttt{sia.status}, \texttt{sia.caller\_method} \\
\hline
\texttt{sia.healthchecks.servers.send} &
Span attorno alla chiamata \texttt{HealthCheckServerToServerApi.SendHealthCheckResult} &
\texttt{sia.operation=send}, \texttt{sia.server\_id}, \texttt{sia.http.status\_code}, \texttt{sia.result}, \texttt{exception.type} \\
\hline
\end{tabularx}
\end{table}

Le trace HTTP generate automaticamente da \texttt{HttpClient} restano figlie di
\texttt{sia.healthchecks.servers.send}, permettendo di ispezionare nella stessa trace sia la
logica applicativa sia il dettaglio del POST verso il \emph{manager}.

\subsection{Metriche}

Il modulo espone tre categorie di metriche: \emph{cumulative} (\texttt{Counter},
\texttt{Histogram}) per \emph{throughput} e durate; \emph{istantanee}
(\texttt{ObservableGauge}) per i valori delle risorse hardware.

\begin{table}[h]
\centering
\caption{Metriche custom del modulo \texttt{Sia.HealthChecks.Servers}}
\label{tab:metrics-healthchecks}
\begin{tabularx}{\textwidth}{|l|l|l|X|}
\hline
\textbf{Metrica} & \textbf{Tipo} & \textbf{Unit} & \textbf{Tag} \\
\hline
\texttt{cycle.count} & Counter & -- & \texttt{sia.success} \\
\hline
\texttt{cycle.duration} & Histogram & \texttt{ms} & \texttt{sia.success} \\
\hline
\texttt{check.count} & Counter & -- & \texttt{sia.check}, \texttt{sia.status} \\
\hline
\texttt{check.duration} & Histogram & \texttt{ms} & \texttt{sia.check}, \texttt{sia.status} \\
\hline
\texttt{check.errors} & Counter & -- & \texttt{sia.check}, \texttt{exception.type} \\
\hline
\texttt{send.count} & Counter & -- & \texttt{sia.result} \\
\hline
\texttt{send.duration} & Histogram & \texttt{ms} & \texttt{sia.result} \\
\hline
\texttt{sql.checks} & Counter & -- & \texttt{sia.connection}, \texttt{sia.result} \\
\hline
\texttt{cpu.usage} & ObservableGauge & \texttt{\%} & -- \\
\hline
\texttt{cpu.temperature} & ObservableGauge & \texttt{Cel} & -- \\
\hline
\texttt{ram.available} & ObservableGauge & \texttt{GBy} & -- \\
\hline
\texttt{ram.total} & ObservableGauge & \texttt{GBy} & -- \\
\hline
\texttt{disk.free} & ObservableGauge & \texttt{GBy} & \texttt{sia.drive} \\
\hline
\texttt{disk.total} & ObservableGauge & \texttt{GBy} & \texttt{sia.drive} \\
\hline
\end{tabularx}
\end{table}

I nomi sono prefissati con \texttt{sia.healthchecks.servers.} (omesso dalla tabella per
compattezza). All'esportazione su Prometheus i punti vengono sostituiti con trattini bassi
e l'\texttt{unit} viene espansa nel nome (es.\ \texttt{sia\_healthchecks\_servers\_cpu\_temperature\_celsius}).

\subsection{Log}

I log prodotti dal \texttt{HostedService} (avvio, iterazione abilitata/disabilitata,
esito invio) sono già propagati via \emph{sink} OTLP di Serilog e correlati alle trace
tramite \texttt{TraceId} e \texttt{SpanId}. Nessuna modifica infrastrutturale è stata
richiesta sul logging; l'intervento si è limitato a sostituire alcune
\emph{string interpolation} con template strutturati per migliorare la ricercabilità su
Loki.

\section{Struttura del modulo}

Il modulo è organizzato attorno alla classe statica
\texttt{SiaHealthChecksServersDiagnostics}, che concretizza il pattern
\texttt{Sia\{Modulo\}Diagnostics} illustrato nella sezione~\ref{sec:design-patterns} e
nella figura~\ref{fig:diag-diagnostics-pattern}. Rispetto al pattern generico, la specializzazione
rilevante per questo modulo è l'utilizzo massiccio degli strumenti \texttt{ObservableGauge}
per esporre i valori \emph{istantanei} delle letture hardware (CPU\%, temperatura, GB,
stato SQL) e il meccanismo di \emph{cache} che disaccoppia la lettura (effettuata dagli
\texttt{IHealthCheck} quando il probe ASP.NET Core chiama \texttt{/health}) dalla
pubblicazione al \emph{Meter} (effettuata on-demand dalla \emph{callback} della \emph{gauge}
al momento dello \emph{scrape} di Prometheus).

I cinque \texttt{IHealthCheck} e il \texttt{HostedService} dipendono dalla classe
\texttt{Diagnostics} per:
\begin{itemize}
    \item creare le \emph{span} tramite l'\texttt{ActivitySource} condiviso (\texttt{cycle},
        \texttt{check}, \texttt{send});
    \item incrementare i \texttt{Counter}/\texttt{Histogram} di esecuzione e durata;
    \item aggiornare la cache dei valori correnti tramite i metodi
        \texttt{UpdateCpuUsage}, \texttt{UpdateRam}, \texttt{UpdateDisk} ecc., che vengono
        poi letti dalle \texttt{ObservableGauge} al prossimo \emph{scrape}.
\end{itemize}
La classe \texttt{HealthCheckServerToServerApi} invocata dal \texttt{HostedService} per
l'invio al \emph{manager} non richiede \emph{span} custom: l'istrumentazione automatica di
\texttt{HttpClient} produce già la \emph{child span} \texttt{http.client} come discendente
della span \texttt{send} creata dal \texttt{HostedService}, preservando la gerarchia
parent/child descritta nel modello delle trace della sezione~\ref{sec:design-patterns}.

\section{Configurazione e attivazione}

Come gli altri moduli, \texttt{Sia.HealthChecks.Servers} è \emph{opt-in} tramite
\emph{appsettings}. Nel pipeline condiviso (\texttt{BaseApplicationBuilder}) è sufficiente
la registrazione del \texttt{Meter} e dell'\texttt{ActivitySource} tramite il
\texttt{ContainerExtensions} del modulo; la configurazione dichiarativa è:

\begin{verbatim}
"OpenTelemetry": {
  "SamplingRatios": {
    "Sia.HealthChecks.Servers": 1.0
  },
  "Instrumentation": {
    "SiaHealthChecksServers": { "Metrics": true }
  }
}
\end{verbatim}

\section{Dashboard Grafana}

La dashboard \texttt{sia-healthchecks-servers.dashboard.json} è versionata insieme al codice
del modulo, nella cartella \texttt{Sia.HealthChecks.Servers/Diagnostics/}. Segue lo schema
standard descritto nel capitolo~\ref{cap:progettazione}, con le \emph{template variable}
obbligatorie \texttt{service\_name}, \texttt{customer\_name}, \texttt{enduser\_id},
\texttt{enduser\_name}, \texttt{bucket}, e i pannelli organizzati nell'ordine:
\emph{overview} $\rightarrow$ durate $\rightarrow$ \emph{throughput} $\rightarrow$ risorse $\rightarrow$ trace $\rightarrow$
log.

La novità rispetto alle altre dashboard di modulo è la riga \emph{Risorse}, dedicata alle
\texttt{ObservableGauge} specifiche del modulo:
\begin{itemize}
    \item andamento CPU\% e temperatura CPU nel tempo;
    \item RAM disponibile e totale con confronto a soglia;
    \item spazio disco libero per drive;
    \item \emph{state timeline} della metrica \texttt{sql.checks} con \emph{break-down} per
        \texttt{sia.connection} e \texttt{sia.result}, per visualizzare a colpo d'occhio gli
        intervalli in cui una connessione SQL è risultata \emph{down}.
\end{itemize}

------------------AGGIUNGERE IMMAGINE DASHBOARD ---------------------------------

I pannelli \emph{trace} e \emph{logs} consentono, a partire da un picco anomalo sulle
metriche \emph{hardware}, di navigare alla trace del ciclo di notifica corrispondente
e ai \emph{log} emessi nell'iterazione, chiudendo il percorso diagnostico
metrica $\rightarrow$ trace $\rightarrow$ log descritto in generale nei
capitoli~\ref{cap:progettazione} e~\ref{cap:codifica}.
